\documentclass[12pt]{jlreq}
\usepackage{amsmath, amssymb}

\newcommand{\C}{\mathbb{C}}
\newcommand{\R}{\mathbb{R}}
\newcommand{\Z}{\mathbb{Z}}
\newcommand{\N}{\mathbb{N}}
\newcommand{\F}{\mathbb{F}}
\newcommand{\Prob}{\mathbb{P}}
\newcommand{\Exp}{\mathbb{E}}
\newcommand{\dd}{\,d}
\newcommand{\Ker}{\operatorname{Ker}}
\newcommand{\Impart}{\operatorname{Im}}
\newcommand{\Hom}{\operatorname{Hom}}

\begin{document}

\((\Omega, \mathcal{B}, \mu)\) を測度空間，
\[
  X = \left\{ f : \Omega \to \R \,\middle|\, f \text{ は } \mathcal{B}\text{-可測で } \int_\Omega \sqrt{|f(x)|} \, \mu(dx) < \infty \text{ を満たす．} \right\}
\]
とする．ただし，\(L^p\)-空間の場合と同様に，ほとんどいたるところ \(f(x) = g(x)\) のときは \(f = g\) とみなす．

(1) \(f, g \in X\) に対して，
  \[
    \rho(f,g) = \int_\Omega \sqrt{|f(x) - g(x)|} \, \mu(dx)
  \]
  と定義すると，\(\rho(\cdot,\cdot)\) は \(X\) 上の距離であることを示せ．

(2) \((X, \rho)\) は，完備な距離空間であることを証明せよ．

\end{document}